Reading the mapping matrix by column rather than by row reveals where the technical
evaluation literature is rich and where it is thin relative to what the frameworks
demand. Table~\ref{tab:gap} summarizes the coverage breadth of evidence for each requirement, and
Fig.~\ref{fig:gap} visualizes the aggregate. We characterize coverage breadth by the
number of method families with \emph{primary} relevance: \emph{broadly covered} requirements
have several such families, \emph{moderately covered} ones one or two, and
\emph{narrowly covered} ones none.

\subsection{Well-served requirements}
The requirements most tightly coupled to the technical literature are those framed as
measurable trustworthiness characteristics: accuracy and performance, robustness,
security, explainability, and fairness. These correspond to NIST \textsc{Measure}~2.x,
EU AI Act Article~15, and the ISO/IEC~42001 verification-and-validation control. The
omnibus technical requirements that draw on several method families at once---EU
Article~15 (accuracy, robustness, and cybersecurity) and ISO/IEC~42001 A.6.2.4
(verification and validation)---reach the highest evidentiary coverage in our scheme
(``broad'' in Table~\ref{tab:gap}), while each individual trustworthiness characteristic
has primary relevance from one or more established method families. For systems built from
well-studied model classes on tabular or vision data, a deployer can today assemble
strong conformity evidence for these obligations using off-the-shelf methods.

\subsection{Process obligations: a definitional boundary, not a methods gap}
A first category is a structural boundary rather than a deficiency in the literature.
Obligations written as \emph{processes}---NIST \textsc{Govern}, EU AI Act Article~9 (the
risk-management system), and the ISO/IEC~42001 management-system clauses---are, by design,
not the kind of thing any single technical evaluation can satisfy. This is not an absence
of methods: documentation (M8, M9) and monitoring (M11) supply substantial evidence
\emph{inputs} to these processes, and the matrix records those supporting and primary
contributions. What process obligations require in addition is an organizational process
that integrates the technical evidence; the ``gap'' is therefore the need for such
integration, not missing measurement. We flag this so practitioners do not mistake a strong
benchmark report for conformity with a process obligation, and we reserve the term
\emph{gap} (below) for requirements that genuinely lack primary-relevance technical methods.

A second, more troubling gap concerns obligations for which a technical evaluation
\emph{should} exist but the methods are immature. \emph{Human-oversight effectiveness}
(EU Article~14) is the clearest case: although a human-in-the-loop is mandated, methods
for evaluating whether oversight is \emph{effective} (whether humans exhibit automation
bias, rely appropriately, or can actually intervene) are nascent~\cite{guo2024reliance,%
chen2024reliancedrills,romeo2025automationbias}. \emph{Record-keeping and logging} (EU
Article~12) is an engineering capability that, in practice, lacks evaluation methods for
assessing whether logs are adequate to reconstruct system behavior. \emph{AI impact
assessment} (ISO/IEC~42001 A.5; EU fundamental-rights impact assessment) has frameworks
but few standardized, quantitative evaluation instruments.

\subsection{The agentic measurement gap}
A significant and rapidly growing gap is created by the shift from models to agents.
Current evaluation methods overwhelmingly grade \emph{outputs}---the text a model emits in
response to a probe---whereas the risks of autonomous, tool-using systems arise from
\emph{actions}: which tools are invoked, which sub-agents are delegated to, how a plan
allocates resources. This is not a speculative concern: independent work shows that a
system near-parity on a question-answering fairness benchmark can still act with disparity
when it selects tools or escalates cases~\cite{li2025actions,blankenstein2025biasbusters},
that agent safety must be assessed over multi-step behavior rather than single
responses~\cite{zhang2024agent}, and that tool-integrated agents open attack surfaces (e.g.,
indirect prompt injection) absent in the base model~\cite{greshake2023notwhat}; dedicated
agent benchmarks are only beginning to appear~\cite{gu2025safesearch,peng2025eu}. Of our
corpus, only a small number of references evaluate action-level behavior, confirming the
class is under-covered. A recent framework (citation withheld for double-blind review) offers
one formalization of this measurement gap (token-level parity coexisting with action-level
disparity), but the gap stands on the independent evidence above. What is novel here is not
the individual agent benchmarks, which we cite, but the systematic identification of this
gap as a \emph{structural} mismatch between frameworks drafted for model-level evaluation
and the deployment reality of agentic, tool-integrated systems, surfaced by reading the
cross-framework matrix by column. Because the EU AI Act,
NIST AI RMF, and ISO/IEC~42001 were drafted with model-centric assessment in mind, their
requirements are presently evidenced by methods that do not observe the behaviors most
likely to cause high-risk harm in agentic deployments. Action-level evaluation harnesses,
trajectory logging tied to the logging obligations, and agentic red-teaming are, on the
evidence above, the most critical near-term directions for aligning technical evaluation
with regulatory intent (Section~\ref{sec:future}).

% Gap-analysis table (Table T5).
\begin{table}[t]
\centering
\caption{Gap analysis: \emph{coverage breadth} of technical evaluation evidence per requirement, with the principal gaps. Coverage breadth counts how many method families supply \emph{primary}-relevance evidence in Table~\ref{tab:matrix}: broad ($\geq$3 families), moderate (1--2), narrow (0). It measures breadth of mapped families, not the methodological maturity or validation quality of the underlying methods. \textbf{Takeaway:} the narrowly-covered requirements are consistently \emph{process} obligations (NIST \textsc{Govern}, EU Article~9, ISO A.5/A.10) and the obligations whose evaluation methods are immature (oversight effectiveness, logging adequacy), not the measurable trustworthiness characteristics.}
\label{tab:gap}
\footnotesize
\begin{tabular}{p{0.10\columnwidth}p{0.30\columnwidth}p{0.16\columnwidth}p{0.30\columnwidth}}
\toprule
FW & Requirement & Coverage breadth & Principal gap\\
\midrule
NIST & NIST GOVERN & Narrow & Governance is process-oriented; technical eval contributes documentation evidence only.\\
NIST & ME 2.10 Priv. & Narrow & Privacy-evaluation methods (membership inference, differential privacy) are established but function as vulnerability tests; full privacy assurance for LLMs/agents is under-served.\\
NIST & ME 2.11 Fair & Moderate & Model-level fairness evaluation is well developed; action-level/agentic disparity measurement is a structural gap.\\
EU & Art.9 & Narrow & EU Article 9 risk management is a process obligation: technical evaluations are inputs to, not a substitute for, the risk-management system.\\
EU & Art.12 & Moderate & Logging/record-keeping is an engineering capability; few evaluation methods assess log adequacy for reconstruction.\\
EU & Art.14 & Moderate & Human-oversight effectiveness evaluation (automation bias, appropriate reliance) is nascent relative to the obligation.\\
ISO & A.5 Impact & Narrow & AI impact-assessment is a process obligation; methods exist but lack standardized, quantitative evaluation instruments, so no family has primary relevance.\\
ISO & A.10 3rd-pty & Narrow & ISO A.10 third-party and customer relationships is served by supplier-assurance processes; few technical evaluation methods are primarily relevant to it.\\
\bottomrule
\end{tabular}
\end{table}

